%%%%%% 开启“外文资料译文”
%% \translatedliterature{<译文标题>}{<原文作者>}
\translatedliterature{关于我的杀父仇人疑似是名震天下的
	大侠时该如何报仇}{杨过}

	%% 这部分内容务必使用模板提供的首字母大写的标题命令生成对应的外文译文标题，
	%% 否则，对应内容将出现在正文的目录和pdf的书签中，非常冗余。
	%% 另外，可使用的各级标题有\Section{}、\Subsection{}、\Subsubsection{}，没有chapter

	\textit{摘要}——杨过少年时期母亲染病而亡

	\textit{关键词}——练武，离经叛道，复仇，抗敌，练武，离经叛道，复仇，抗敌，练武，离经叛道，复仇，抗敌，练武，离经叛道，复仇，抗敌

	\Section{引言}

	\begin{table}[htbp]
		\captionsetup{list=no}% 阻止此表格显示在表目录中
		\caption{表格样例}
		\begin{tabular}{ccccc}
			\toprule
			\multirow{2}{*}{Column0} &  \multicolumn{2}{c}{Column1\tnote{1}} & \multicolumn{2}{c}{Column2\tnote{2}} \\
			\cmidrule(lr){2-3}\cmidrule(l){4-5}
			~     & subcolumn1 & subcolumn2 & subcolumn1 & subcolumn2 \\
			\midrule
			Row1  & element11 & element12 &element13 & element14 \\
			Row2  & element21 & element22 &element23 & element24 \\
			\cmidrule{2-3}\cmidrule{4-5}
			Row3  & element31 & element32 &element33 & element34 \\
			\bottomrule
		\end{tabular}
	\end{table}


	\Subsection{贡献}


	\Subsubsection{好好好}

	\begin{longtable}{p{2em} p{4.5em}}
		\captionsetup{list=no}% 阻止此表格显示在表目录中
		\caption{中科院部分推荐期刊及会议}\\
		
		\toprule
		\textbf{序号} & \textbf{简称} \\
		\midrule
		\endfirsthead
		
		% 在这里设计首页以外的表题和表头
		\CPcaption{2}{中科院部分推荐期刊及会议}\\
		\toprule
		\textbf{序号} & \textbf{简称} \\
		\midrule
		\endhead
		
		% 在这里设计首页以外的表尾
		\bottomrule
		\multicolumn{2}{l}{续下页} \\  % 如不希望跨页表尾显示任何内容则注释掉即可
		\endfoot
		
		\bottomrule
		\endlastfoot
		
		1 & JSAC \\
		2 & TMC \\
		3 & TON \\
	\end{longtable}
	
	\Section{相关工作}

	\begin{equation}
		x^2 + y^2 + z^2
	\end{equation}
	
	\Section{系统模型}

	\Section{算法设计}

	\begin{algo}[!h]
		\caption{algo环境伪码示例}
		\Input{1) 输入1; 2) 输入2。}
		\Output{输出结果。}

		伪码行1。
		
		\For(\tcc*[f]{循环条件注释1}){循环条件1}{
			伪码行2。
			
			\tcp{注释2}
			伪码行3。
			
			\DoWhile(\tcc*[f]{循环条件注释3}){循环条件2}{
				伪码行4。
			}
			
			\tcc{loop循环}
			\Loop(\tcc*[f]{注释4}){
				循环体1。
			}
			
			\Repeat(\tcc*[f]{循环条件注释5}){循环条件3}{
				循环体2。
			}
			\eIf(\tcc*[f]{条件注释6}){条件语句6}{
				为真，伪码行5。
			}{
				条件为假，伪码行6。\tcp*[f]{repeat循环}
			}
		}
		\textbf{return} 算法结果。
	\end{algo}
	
	\Section{实验}


	
	\Section{结论}